\section*{Abstract}
\addcontentsline{toc}{section}{Abstract}

The HIBEAM/NNBAR programme is developing an annihilation detector for neutron-conversion searches, including free neutron--antineutron oscillation. One proposed detector element is a range stack of plastic-scintillator staves intended to sample the penetration depth and longitudinal energy-loss pattern of low-energy charged particles. This paper studies the CCB test-beam range-stack prototype using beam data and Geant4-based detector simulations.

The analysed B-stack readout uses channels labelled B2, B4, B6 and B8, while the corresponding detector model contains eight physical longitudinal layers. Existing source-bound selected-data products indicate different longitudinal occupancy patterns for the declared Sample-I and Sample-II run families, with one population concentrated closer to the stack entrance and the other extending more frequently to downstream readout positions. The final quantitative depth-profile result is withheld until regeneration from the pre-threshold 8-channel-by-16-sample event product with run-aware uncertainty and closed polarity/mapping contracts. Primary run-log and hardware-trigger records are also incomplete, so any surviving difference must initially be interpreted as a run-family observation rather than a measured trigger efficiency.

The physically relevant amplitude telescope observable is \(\Delta E=A(B2)\) and \(E=A(B4)+A(B6)+A(B8)\). Sparse longitudinal readout can hide the layer containing the largest energy loss and therefore distort a stopping/range proxy. The current numerical \(\Delta E\)--\(E\) data/MC bundle is under publication hold because the beam event construction is threshold-censored upstream and the MC producer does not preserve an unambiguous physical-layer namespace or validated event-measure contract.

The located beam waveform product contains eight channels with 16 samples per channel. A direct B4--B6 analysis of the current reconstructed subset gives a central residual width of approximately 38 ns. This is reported only as a waveform-format/estimator-level pair residual: component identity, timing reference, run dependence and covariance do not support conversion to an intrinsic single-stave resolution.

A detailed single-stave optical model provides a framework for separating quenching, scintillation, wavelength shifting, optical transport, SiPM response and electronics. The historical July optical grid and its derived held-out energy-reconstruction result are superseded model diagnostics because the campaign predates source-level optical repairs and because its target variable was Birks-visible energy rather than unquenched deposited energy. Absolute optical response and detector energy resolution are therefore withheld pending regeneration of the current model, explicit response nuisances and a genuinely held-out reconstruction study.
